\begin{abstract}
Only 1 in 1,000 Python files in a curated code corpus contains an independently-executable
doctest --- a 31$\times$ gap below the 3.1\% rate at which \texttt{>>>} doctest patterns appear.
We document this finding through a systematic three-phase feasibility scan of 10,000 randomly
sampled Python files: Phase A (pattern detection: 3.1\%), Phase B (AST extraction: 2.0\%), and
Phase C (subprocess execution: 0.1\%). The collapse at Phase C is driven by import isolation:
Python library code universally depends on third-party packages unavailable in clean subprocess
environments. The resulting executable-doctest token pool (0.004M tokens) is 125,000$\times$
below the 500M-token SFT budget target, making doctest-passing filtering structurally infeasible
for raw Python corpus SFT construction. We validate \texttt{compile()}-based filtering as the
practical alternative --- scalable to 12.96M+ files, requiring no dependency management, and
providing a reproducible syntax-validity quality gate. We release the validated three-phase
pipeline (28/28 tests, 129.8s/10K files) and design the equal-token-budget SFT comparison
experiment (\texttt{compile}-only vs.\ unfiltered, Qwen2.5-Coder-1.5B, HumanEval/MBPP pass@1)
as immediate future work.
\end{abstract}
